\chapter{1923:Fritz Pregl}
\label{chap:chemistry1923}

The Nobel Prize in Chemistry for the year 1923 was awarded to Fritz Pregl.

\textbf{Motivation:}
for his invention of the method of micro-analysis of organic substances

\section{Fritz Pregl}

\subsection*{Early Life and Education}

Fritz Pregl was born on 1869-09-03 in Laibach, Austria-Hungary (now Ljubljana, Slovenia).
Fritz Pregl, was a Slovenian-Austrian chemist and physician from a mixed Slovene-German-speaking background. He won the Nobel Prize in Chemistry in 1923 for making important contributions to quantitative organic microanalysis, one of which was the improvement of the combustion train technique for elemental analysis.

\subsection*{Career and Major Research}

Pregl was professor of medical chemistry at the University of Graz. He perfected quantitative micro-analysis, reducing the sample size required for elemental analysis of organic compounds to a few milligrams. This technique made possible the chemical analysis of the small quantities of substances that occur in biological materials.

\subsection*{Nobel Prize-winning Work}

The work for which Fritz Pregl was awarded the Nobel Prize in Chemistry in 1923 was described by the Nobel Committee as: "for his invention of the method of micro-analysis of organic substances".

This research fundamentally advanced the field by making groundbreaking contributions that transformed chemical science.

\subsection*{Significance and Impact}

The impact of Fritz Pregl's work extends far beyond the specific achievement recognized by the Nobel Prize.
These contributions continue to influence chemical research and education today.

\subsection*{Later Career and Honors}

Fritz Pregl continued his scientific work until 1930-12-13, passing away in Graz, Austria.
He received numerous honors including membership in major scientific academies and societies.

\subsection*{Legacy}

The legacy of Fritz Pregl endures through the lasting impact of his scientific contributions.
Fritz Pregl's work continues to be cited and built upon by researchers across the chemical sciences.

\subsection*{Modern Developments}

Pregl's microanalysis, which scaled elemental analysis down to milligram samples, transformed analytical chemistry and enabled the characterization of scarce natural products and biologically important molecules. His techniques evolved into modern microanalysis, combustion analysis, and the automated elemental analyzers used in pharmaceutical and materials laboratories today.

\subsection*{Selected Publications}

Fritz Pregl's major publications include numerous papers in leading journals of the era, detailing the experimental and theoretical work summarized above.

\clearpage

\section*{Chapter Summary}

Fritz Pregl was awarded the prize for his invention of the method of micro-analysis of organic substances.
